% vim: spl=pt
\begin{exercício}{Espectro de operadores auto-adjuntos}{ex21}
    Seja \(T \in \bounded(\hilbert)\) um operador auto-adjunto. Mostre que
    \begin{enumerate}[label=(\alph*)]
        \item \(\sigma(T) \subset \mathbb{R}\) e
        \item \(\sigma_r(T) = \emptyset.\)
    \end{enumerate}
\end{exercício}
\begin{proof}[Resolução de (a)]
  Seja \(\lambda \in \mathbb{C} \setminus \set{\mathbb{R}},\) então para todo \(x \in \hilbert\) com \(\norm{x} = 1\) vale
  \begin{equation*}
    \abs{\inner{x}{(\lambda \unity - T)x}} \leq \norm{x} \norm{(\lambda \unity - T)x} = \norm{(\lambda \unity - T)x},
  \end{equation*}
  portanto
  \begin{equation*}
    \norm{(\lambda \unity - T)x} \geq \abs{\inner{x}{(\lambda \unity - T)x}} = \abs{\lambda - \inner{x}{Tx}} \geq \abs{\Im\lambda}.
  \end{equation*}
  Assim, mostramos que \(\lambda \unity - T\) é limitado inferiormente e, portanto, é injetivo e com imagem fechada. Como \(T\) é auto-adjunto, \(\conj{\lambda} \unity - T\) também é injetivo, portanto
  \begin{equation*}
    \ran(\lambda \unity - T)^\perp = \ker(\conj{\lambda} \unity - T) = \set{0},
  \end{equation*}
  o que nos informa que a imagem de \(\lambda \unity - T\) é densa em \(\hilbert.\) Como a imagem de \(\lambda \unity - T\) é fechada, este operador é sobrejetor. Ainda, como \(\lambda \unity - T\) é limitado inferiormente, a sua inversa é contínua, e concluímos que \(\lambda \notin \sigma(T).\) Isto é, \(\mathbb{C} \setminus \set{\mathbb{R}} \subset \mathbb{C} \setminus \sigma(T),\) o que conclui a demonstração.
\end{proof}
\begin{proof}[Resolução de (b)]
  Para um operador limitado \(A \in \bounded(\hilbert),\) vale \(\lambda \in \sigma_r(A) \implies \conj{\lambda} \in \sigma_p(A^*),\) portanto o complexo conjugado de \(\sigma_r(A)\) está contido em \(\sigma_p(A^*).\) Para o operador auto-adjunto e limitado \(T,\) temos \(\sigma_r(T) \subset \sigma_p(T),\) mas \(\sigma_r(T) = \sigma_r(T) \cap \sigma_p(T) = \emptyset,\) como desejado.
\end{proof}
